% =============================================================================
% DATASET + CONFOUNDERS   --   target ~1.25 pp. Borrowed from paper/ and cut.
% Cuts vs paper/ (logged in docs/faimi/dataset-section-cuts.md): dropped
% tab:notation (symbols inlined), fig:volumes, tab:volumes (folded to one
% sentence), the connected-components paragraph, the GE/Siemens voxel arithmetic
% (-> 1 line). KEPT: fig:anatomy, tab:cohort. Added: gold/silver comparability
% (load-bearing for E2) -- still needs an empirical one-liner.
% =============================================================================
\section{Dataset and Exploratory Analysis}
\label{sec:dataset}

\subsection{The CSpineSeg Cohort}
\label{sec:cohort}

We work with CSpineSeg~\cite{zhou2025cspineseg}, a public collection of 1,255
sagittal T2-weighted cervical-spine MRI exams from 1,232 patients at Duke
University. We drop a single localizer series with anomalous dimensions, leaving a
working set $\mathcal{D}_0$ of $N_0 = 1{,}254$ exams from 1,231 patients. The task
is semantic segmentation: labelling every voxel of a scan as background, vertebral
body, or intervertebral disc ($\mathcal{L} = \{0, 1, 2\}$; \autoref{fig:anatomy}).
About $40\%$ of the exams carry expert (gold) reference segmentations, drafted by one post-doctoral researcher without medical training and reviewed and revised by
one of six board-certified radiologists; the rest were segmented automatically by the authors' own model, and
these silver masks were never human-reviewed -- though the mid-sagittal slice of
each unannotated scan was checked for image quality.

We write the cohort as $\mathcal{D} = \{(X_i, Y_i, A_i, q_i)\}_{i=1}^{N}$: exam
$i$ pairs a 3D image $X_i$ with a reference segmentation $Y_i$, an attribute value
$A_i$ over the protected attributes $\mathcal{A} = \{\text{sex}, \text{race},
\text{age}\}$, and a provenance flag $q_i \in \{\mathsf{g}, \mathsf{s}\}$ marking a
gold (expert) or silver (machine-generated) label. We write
$\mathcal{D}_{A=a}$ for the subgroup with $A = a$. The analysis cohort
$\mathcal{D}$ ($N = 1{,}142$) is the sex-balanced subset of $\mathcal{D}_0$ defined
in ``\nameref{sec:regimes}''.

\begin{figure}[t]
\centering
\includegraphics[width=0.55\textwidth]{anatomy_overlay.png}
\caption{A representative mid-sagittal T2 slice from CSpineSeg with the expert
reference segmentation overlaid: vertebral bodies (red, class~$1$) and
intervertebral discs (green, class~$2$); every other voxel is background. Each
exam is a 3D volume; a single slice is shown for illustration.}
\label{fig:anatomy}
\end{figure}

\autoref{tab:cohort} summarises $\mathcal{D}_0$. The cohort leans slightly female
(55\%), is predominantly White (65\%) with a sizeable Black minority (28\%) and
only small numbers in the remaining race groups, and is overwhelmingly
non-Hispanic. Ages span the adult range (18--89, median 56), and scans are split
across two manufacturers (Siemens 63\%, GE 37\%) and two field strengths (1.5\,T
60\%, 3.0\,T 40\%). The volumes are strongly anisotropic: in-plane resolution is
high (${\sim}0.53$\,mm) but the sagittal slices are thick (${\sim}4$\,mm) and few
(12--25 per scan). These group sizes decide which fairness axes we can test: sex
(Female vs.\ Male), race (White vs.\ Black), and age have enough exams per group
for reliable inference, so we test them formally; the smaller race categories
(8\% combined) and ethnicity are too sparse and are reported descriptively only.

\begin{table}[t]
\centering
\caption{Composition of the working set $\mathcal{D}_0$ ($N_0 = 1{,}254$ exams,
counted at the exam level). Percentages are of $N_0$.}
\label{tab:cohort}
\small
\begin{tabular*}{\textwidth}{@{\extracolsep{\fill}} l r l r}
\hline\noalign{\smallskip}
Characteristic & $n$ (\%) & Characteristic & $n$ (\%) \\
\noalign{\smallskip}\hline\noalign{\smallskip}
\textit{Sex} & & \textit{Age, years}$^{b}$ & $54.6 \pm 16.3$ \\
\quad Female & 683 (54.5) & \quad $<$40 & 241 (19.2) \\
\quad Male & 571 (45.5) & \quad 40--60 & 486 (38.8) \\
\textit{Race} & & \quad $\ge$60 & 527 (42.0) \\
\quad White & 809 (64.5) & \textit{Manufacturer} & \\
\quad Black & 349 (27.8) & \quad Siemens & 788 (62.8) \\
\quad Other/unknown$^{a}$ & 96 (7.7) & \quad GE & 466 (37.2) \\
\textit{Ethnicity} & & \textit{Field strength} & \\
\quad Non-Hispanic & 1{,}162 (92.7) & \quad 1.5\,T & 746 (59.5) \\
\quad Hispanic & 59 (4.7) & \quad 3.0\,T & 508 (40.5) \\
\quad Not reported & 33 (2.6) & & \\
\noalign{\smallskip}\hline
\end{tabular*}

\smallskip
{\footnotesize $^{a}$\,Asian (25), American Indian or Alaska Native (9), Native
Hawaiian or other Pacific Islander (1), Other (28), Not Reported (33).
\quad $^{b}$\,mean $\pm$ SD; 13 exams with missing age (confirmed ${>}89$) are
grouped into the oldest bin.}
\end{table}

\subsection{Anatomical Confounders}
\label{sec:confounders}

A performance gap is bias only if it is not explained by how the groups differ
anatomically or in acquisition. Anatomy can make a group genuinely harder --
cervical degeneration rises with age~\cite{kato2012cervical}, bridging and merging
adjacent vertebrae and blurring the boundaries the model must trace -- and can also
fool the metric, since Dice rewards large structures (a fixed boundary error costs
a smaller fraction of a big structure's score than a small one's). Acquisition
matters only if a group were scanned mostly on particular hardware; imaging site, a
confounder for multi-centre studies, does not arise in single-institution
CSpineSeg.

We measured the physical volume (mm$^3$) of the vertebral bodies and discs in
every $\mathcal{D}_0$ reference and compared across $\mathcal{A}$, reading
rank-based effect sizes ($r_{rb}$~\cite{wendt1972rankbiserial} for sex/race,
$\varepsilon^2$~\cite{tomczak2014effectsize} for age) on the conventional
negligible/small/medium/large scale~\cite{cohen1988power}. Sex drives a large gap
in both structures (men ${\sim}23$--$24\%$ larger); race and age enlarge the
vertebral body only (White $+12\%$, small; $\ge$60 vs.\ $<$40 $+20\%$, medium)
while leaving disc volume essentially unchanged -- making the disc the structure
least exposed to the size confound. A quick check rules out the scanner: GE yields ${\sim}2.15\times$
Siemens' voxel count purely from finer in-plane resolution, but in physical mm$^3$
the volumes align across manufacturers (${\sim}6\%$ gap), so these differences are
real anatomy, not a resolution artefact.

No attribute is otherwise tangled up with the scanner, which could fake a group
gap: race is independent of both manufacturer and field strength, and sex of race
(all Cram\'er's $V \approx 0.05$--$0.07$, none significant). The one mild
entanglement is age -- White patients are slightly older than Black (median $58$
vs.\ $54$) -- so because age tracks vertebral-body anatomy, we control for age
whenever we compare race. Importantly for the biased-ruler experiment
(``\nameref{sec:design}''), the gold and silver pools are themselves
demographically comparable: Zhou et al.~chose which exams to expert-annotate
pseudo-randomly, by medical-record number rather than any demographic or
difficulty criterion~\cite{zhou2025cspineseg}, so a gold-vs-silver performance gap
reflects label provenance, not a difference in case mix.
% TODO: back the comparability claim with one empirical line -- gold vs silver
% pools match on sex/race/age marginals (and ideally a difficulty proxy). This
% pre-empts the "maybe gold cases were just harder" objection that the ~8-pt E2
% gap will otherwise invite. See docs/faimi/paper-framing.md (risk section).

Two choices follow and carry into the rest of the paper. First, we report the
vertebral body and disc separately as well as their macro average; because the
disc is anatomically uniform across race and age, it is the structure least
confounded by size and so serves as a cleaner cross-check on any macro-level gap.
Second, we evaluate with a volume-aware score -- the normalised Dice score (nDSC,
``\nameref{sec:metrics}'') -- alongside raw Dice, to neutralise the size confound
before any performance gap can be read as bias.
